% Complete finite proofs for P16-C1.
\section{Finite logical abduction: proofs}
\label{sec:annex-p16-finite-abduction}

Use the definitions of Chapter~\ref{sec:p16-finite-abduction}. The source
model and its single-deletion criterion are Eiter and Gottlob's
\cite[sec.~2, manuscript pp.~9--11]{EiterGottlobAbductionManuscript}.
The finite enumeration, extended inconsistent-input case and examples are
proved explicitly here.

\begin{annexitem}{annex:p16.abduction-semantics}{thm:p16-abduction-semantics}{Finite valuations and the explanation criterion}
\begin{proof}
\textbf{Agreement with the book's semantics.}\enspace
Let \(v:V\to\mathbb{B}\). Extend it to a valuation
\(\widetilde v:\mathsf{Prop}\to\mathbb{B}\) by giving every atom outside \(V\)
the value \(\mathbf{F}\). For each formula over \(V\), its recursively
computed value under \(v\) equals its value under \(\widetilde v\).
This follows by induction on the formula: the atomic values agree,
\(\bottom\) has the same fixed value, and each connective applies the
same truth function to the values of its immediate subformulas. Conversely,
restriction of any valuation on \(\mathsf{Prop}\) to \(V\) preserves the
values of all formulas over \(V\) by the same induction. This is also
the finite restriction instance of \cref{lem:prop-coincidence}.

It follows that \(\Gamma\) has a global satisfying valuation exactly
when \(\operatorname{Mod}_V(\Gamma)\) is nonempty. For a formula
\(\varphi\) over \(V\), \(\Gamma\entails\varphi\) holds
exactly when every member of \(\operatorname{Mod}_V(\Gamma)\)
satisfies \(\varphi\): extend a local countervaluation for one
direction and restrict a global countervaluation for the other. Apply
this to \(\Gamma=T\cup E\) and to each atom \(m\in M\).
An atom has value \(\mathbf{T}\) exactly when the valuation satisfies it,
which proves the finite criterion.

The source's theory is finite and consistent, its \(H,M\) are subsets
of the same finite variable alphabet, and its solution condition is exactly
consistency of \(T\cup E\) and consequence of every manifestation.
Its finite alphabet can be named by distinct members of \(\mathsf{Prop}\);
transporting valuations along that naming preserves atomic truth values.
The source's truth constant is translated to \(\neg\bottom\), which
is always true. Its biconditional is translated to the conjunction of the
two material implications. That conjunction is true exactly when the two
formula values agree: equal values make both implications true, and unequal
values make the implication from the true value to the false value false.
Induction on the source formula now proves preservation of truth under
these expansions, since the other constructors have the same truth clauses.
The naming sends subsets of hypotheses bijectively to subsets of the named
hypotheses and preserves both inclusion and cardinality. Thus the criterion
agrees on every source input under this stated notation translation.
The only additional input
case introduced in this chapter drops the consistency requirement on \(T\);
the next paragraph determines that case directly.

\paragraph{Boundary cases.}
Every model of \(T\cup E\) is a model of \(T\), so
\[
 \operatorname{Mod}_V(T\cup E)\subseteq\operatorname{Mod}_V(T).
\]
If the latter is empty, so is the former, for every candidate. Thus an
inconsistent background permits no explanation. In this case semantic
entailment is vacuous, but the distinct consistency conjunct is false.
Putting \(E=\varnothing\) gives \(T\cup E=T\), which proves the
empty-explanation equivalence. If \(M=\varnothing\), every valuation
satisfies every member of \(M\), because it has no members; explanation
then reduces precisely to consistency of \(T\cup E\). All arguments
also apply when \(V\) is empty: there is exactly one function from the
empty set to \(\mathbb{B}\).
\end{proof}
\end{annexitem}

\begin{annexitem}{annex:p16.abduction-minimality}{thm:p16-abduction-minimality}{Minimal explanations and single deletions}
\begin{proof}
\textbf{Qualified monotonicity.}\enspace
If \(E\subseteq F\), every valuation satisfying \(T\cup F\)
satisfies \(T\cup E\). Therefore, if \(E\) explains the observations,
every model of \(T\cup F\) satisfies \(M\). When \(T\cup F\)
is also consistent, this proves that \(F\) is an explanation. The
consistency condition cannot be inferred from \(E\subseteq F\).

\paragraph{Finite minimality.}
Fix \(E\in\operatorname{Expl}(\mathcal P)\). The family
\[
 \{F\subseteq E:F\in\operatorname{Expl}(\mathcal P)\}
\]
is finite and nonempty, since it contains \(E\). Select a member
\(E_0\) of least cardinality in this family. A proper subset of
\(E_0\) which were an explanation would also be a subset of \(E\)
and would have smaller cardinality. That contradicts the choice. Hence
\(E_0\) is subset-minimal among all explanations and is contained in
\(E\).

If the full explanation family is nonempty, it too is a finite subfamily
of \(2^H\), so it has a member of least cardinality. Every such member
is subset-minimal, because a proper explanatory subset would have smaller
cardinality. This proves existence of both families and the asserted
inclusions. If two distinct subset-minimal explanations were comparable,
one would be a proper explanatory subset of the other, contradicting the
latter's minimality. Their cardinalities need not agree: the complete
calculation for \(\mathcal P_{\rm sizes}\) below gives sizes one and two.
That calculation uses the definitions directly, without invoking this
incomparability assertion.

\paragraph{Single-deletion characterization.}
Let \(E\) be an explanation. If \(E\) is subset-minimal, every
\(E\setminus\{h\}\), for \(h\in E\), is a proper subset and
cannot be an explanation.

Conversely, suppose \(E\) is not subset-minimal. Choose an explanation
\(F\subsetneq E\) and some \(h\in E\setminus F\). Put
\(D=E\setminus\{h\}\). Then \(F\subseteq D\subsetneq E\).
A model of \(T\cup E\) exists and also satisfies \(T\cup D\),
so \(T\cup D\) is consistent. Every model of \(T\cup D\)
satisfies \(T\cup F\), and hence satisfies \(M\), because \(F\)
is an explanation. Thus \(D\) is an explanation, contradicting failure
of all one-atom deletions. This proves the converse, including the
consistency step in the source's criterion. If \(E=\varnothing\),
there is no proper subset and no atom to delete; both conditions hold.
If the background is inconsistent, there is no explanation \(E\) to
which the characterization would apply, so the extended convention adds
no exception to the statement.

Finally, the constraint example below gives an explanatory \(\{a\}\)
whose permitted superset \(\{a,b\}\) is inconsistent. It demonstrates
the stated limitation of unqualified upward closure by a direct calculation.
\end{proof}
\end{annexitem}

\begin{annexitem}{annex:p16.abduction-enumeration}{thm:p16-abduction-enumeration}{Soundness, completeness and finite termination of the construction}
\begin{proof}
Let the fixed lists of hypotheses and variables have lengths \(r\) and
\(n\), respectively. Each binary vector of length \(r\) picks a
unique subset of \(H\), and every subset has exactly its membership
vector. Each binary vector of length \(n\) likewise determines a unique
valuation on \(V\), and every valuation has exactly that vector. A
length-zero vector is unique, so these correspondences include the empty
cases. Induction on the length shows that the numbers of vectors are
\(2^r\) and \(2^n\): each extension by one position has the two
choices zero and one, starting from one empty vector.

For a fixed candidate \(E\), induction on the number of scanned
valuations gives two flag invariants. The first flag is true exactly when
one of the scanned valuations satisfies \(T\cup E\). The second is
true exactly when one of them satisfies \(T\cup E\) and falsifies
at least one \(m\in M\). Both invariants hold before any scan because
the flags are false and no valuation has been scanned. At the next
valuation, the two stated tests add exactly its corresponding witness, and
the flags retain all previous witnesses. This proves the induction step.

After the complete valuation scan, the first flag is true exactly when
\(\operatorname{Mod}_V(T\cup E)\) is nonempty, and the second is
false exactly when no model in this set falsifies an observation. By
\cref{thm:p16-abduction-semantics}, their tested combination holds exactly
when \(E\) is an explanation. Thus every appended set is an explanation
(soundness), and every explanation is appended when its unique candidate
vector is processed (completeness). No candidate is processed twice, and
each candidate is appended at most once. Hence
\(\mathcal L=\operatorname{Expl}(\mathcal P)\) with no duplicates.

The subset filter keeps an explanation \(E\) exactly when the complete
list contains no explanatory proper subset of it. Since the list is
complete, this is precisely the definition of
\(\operatorname{MinExpl}(\mathcal P)\); both directions follow from
that equivalence. If \(\mathcal L\) is empty, there is no cardinality
minimum to choose and the stated empty output agrees with its definition.
Otherwise a finite scan maintains the least cardinality seen so far,
starting with the cardinality of the first member. Comparing each
successive cardinality with this value preserves that invariant, so the
final value is the least over the entire list. Retaining exactly the
members of that size gives every minimum-cardinality explanation and only
those explanations. Neither filtering step duplicates members.

Every formula of \(T\) is a finite tree, so recursive truth evaluation
terminates by induction on that tree. The theory, candidate and observation
sets are finite, so the tests at a candidate--valuation pair terminate.
There are exactly \(2^r\cdot2^n=2^{r+n}\) such pairs in the prescribed
unshortened scans. Each of the two filtering procedures scans finite lists
and finite sets. Consequently the whole construction terminates, including
inconsistent theories and empty inputs. This proves a bound on the stated
pair inspections, without replacing formula evaluation or set comparison
by unit-cost assumptions.
\end{proof}
\end{annexitem}

\begin{annexitem}{annex:p16.abduction-examples}{thm:p16-abduction-examples}{Complete calculations for all six inputs}
\begin{proof}
\textbf{Inconsistent background.}\enspace
For \(\mathcal P_{\rm inconsistent}\), no valuation satisfies both
\(a\) and \(\neg a\). Thus neither candidate, \(\varnothing\)
or \(\{a\}\), is consistent with \(T\). There is no explanation
and both selected subfamilies are empty. For each candidate the entailment
of \(m\) alone is true vacuously, since its premise set has no model.

\paragraph{Observation already given.}
The only candidates for \(\mathcal P_{\rm given}\) are
\(\varnothing\) and \(\{a\}\). Both are consistent with
\(T=\{m\}\): assign \(m\) true and, when needed, \(a\) true.
Every model of either augmented theory makes \(m\) true because it is
already in \(T\). Thus both candidates explain. The empty one is a
proper subset of \(\{a\}\) and has size zero, so it is the only
member of both minimality families.

\paragraph{Empty data.}
For \(\mathcal P_{\rm empty}\), the only candidate is the empty set,
and the only valuation on \(V\) is the empty function. It satisfies
the empty theory and every member of the empty observation set. Thus the
empty candidate is an explanation. Having no competitors, it belongs to
both minimality families. Each of the three families therefore has exactly
one member, the empty set.

\paragraph{Unequal sizes.}
Consider any \(E\subseteq\{a,b,c\}\). Assigning \(a,b,c,m\) all
true satisfies \(T\cup E\), so every candidate is consistent. If
\(a\in E\), every model of \(T\cup E\) must make \(m\) true
by \(a\to m\). If \(b,c\in E\), the same conclusion follows
from \((b\land c)\to m\).

For the converse, suppose \(a\notin E\) and
\(\{b,c\}\nsubseteq E\). Assign an atom among \(a,b,c\) true
exactly when it belongs to \(E\), and assign \(m\) false. Then
\(a\) is false and at least one of \(b,c\) is false. Both antecedents
in \(T\) are false, so both implications hold. This valuation satisfies
\(T\cup E\) but falsifies \(m\). Hence \(E\) is not an
explanation. We have proved the exact criterion
\[
 E\in\operatorname{Expl}(\mathcal P_{\rm sizes})
   \quad\Longleftrightarrow\quad
       a\in E\ \text{ or }\ \{b,c\}\subseteq E.
\]
The four subsets containing \(a\) and the one remaining subset
\(\{b,c\}\) are exactly the five listed explanations. The empty set
does not explain, so \(\{a\}\) is subset-minimal. Every proper
subset of \(\{b,c\}\) omits \(a\) and at least one of \(b,c\),
so \(\{b,c\}\) is also subset-minimal. Each other listed explanation
properly contains \(\{a\}\). Thus there are exactly two subset-minimal
explanations, and neither contains the other. Their sizes are one and two,
so the unique minimum-cardinality explanation is \(\{a\}\).

\paragraph{Consistency constraint.}
The four candidates for \(\mathcal P_{\rm constraint}\) are
\[
 \varnothing,\{a\},\{b\},\{a,b\}.
\]
The candidate
\(\{a\}\) is consistent: choose \(a,m\) true and \(b\) false.
It entails \(m\) with \(T\) by the implication \(a\to m\).
For the empty candidate choose \(a,b,m\) all false; this is a model
of \(T\) which falsifies \(m\). For \(\{b\}\), choose
\(b\) true and \(a,m\) false; this is a model of \(T\cup\{b\}\)
which falsifies \(m\). These candidates are consistent but do not
explain. No valuation satisfies \(a,b\) and \(\neg(a\land b)\),
so \(\{a,b\}\) is inconsistent with \(T\). Therefore
\(\{a\}\) is the sole explanation and the sole member of both
minimality families. Its failed superset verifies the consistency qualification.

\paragraph{Permitted hypotheses.}
The only permitted candidates for \(\mathcal P_{\rm permitted}\) are
\(\varnothing\) and \(\{a\}\). In each case set \(z,m\) false
and set \(a\) true if the candidate requires it. The implication
\(z\to m\) holds, the candidate is satisfied, and \(m\) is false.
Thus both candidates are consistent and non-explanatory, and all three
explanation families are empty.

Separately, setting \(z,m\) true provides a model of \(T\cup\{z\}\),
and every model of that set makes \(m\) true by \(z\to m\).
Nevertheless \(\{z\}\nsubseteq H=\{a\}\). It is therefore not
an explanation for the declared input. Allowing it would require a different
hypothesis set and hence a different problem.
\end{proof}
\end{annexitem}

\begin{annexitem}{annex:p16.abduction-truth-boundary}{thm:p16-abduction-truth-boundary}{A countervaluation to recovering hypotheses from observations}
\begin{proof}
The preceding calculation proves that \(\{a\}\) and \(\{b,c\}\)
are subset-minimal explanations in \(\mathcal P_{\rm sizes}\).
Define a valuation on its four atoms by
\[
 v(a)=v(b)=v(c)=\mathbf{F},\qquad v(m)=\mathbf{T}.
\]
Both antecedents \(a\) and \(b\land c\) are false, so the two
implications in \(T\) are true. The observation \(m\) is also true.
Thus \(v\) satisfies \(T\cup M\), while it falsifies both
\(a\) and \(b\land c\). Extending this valuation to \(\mathsf{Prop}\)
as in \cref{thm:p16-abduction-semantics} supplies countervaluations in
the book's global semantics. This proves both failures of entailment.
\end{proof}
\end{annexitem}
